\section{Các Luồng Jenkins CD}

Đồ án có ba luồng CD chính: dev auto deploy, staging release có duyệt, và
developer preview on-demand. Điểm chung là tất cả đều đi qua GitOps commit; không
có job nào trực tiếp apply manifest vào namespace do ArgoCD quản lý.

\subsection{Dev Auto Deploy}

Khi app repo \texttt{main} thay đổi, Jenkins:

\begin{enumerate}
    \item Detect modules thay đổi và chạy CI gates.
    \item Build/push Docker image cho service deployable với tag short commit SHA.
    \item Push thêm tag tiện dụng \texttt{main} và \texttt{latest}.
    \item Clone CD repo, chạy \path{scripts/update-image-tag.sh} \texttt{dev <service> <sha>}.
    \item Commit \texttt{cd(lab2): update dev image tags [skip ci]} vào \texttt{yas-cd/main}.
    \item ArgoCD \texttt{yas-dev} tự sync.
\end{enumerate}

Việc overlay dev dùng commit SHA thay vì \texttt{latest} giúp mỗi deploy có dấu vết
rõ ràng và tránh tình trạng Kubernetes không rollout vì manifest không đổi.

\begin{figure}[H]
\centering
\reportimg{img/jenkins/main_build_success.png}
\caption{Jenkins main build chạy qua CI gates, Docker push và GitOps update}
\end{figure}

\begin{figure}[H]
\centering
\reportimg{img/dockerhub/dev_commit_sha_tags.png}
\caption{Docker Hub hiển thị image tag theo commit SHA cho dev}
\end{figure}

\subsection{Staging Release}

Staging chỉ nhận release tag \texttt{vX.Y.Z}. Jenkins release tag build là
promote-only:

\begin{enumerate}
    \item Tag \texttt{vX.Y.Z} trigger Jenkins multibranch/tag job.
    \item Pipeline bỏ qua test/scan/build lại source vì source đã được gate-check bởi main build.
    \item Jenkins đọc tag nguồn từng service từ \texttt{overlays/dev}.
    \item Dùng \texttt{docker buildx imagetools create} retag image hiện có sang \texttt{vX.Y.Z}.
    \item Chạy \path{scripts/promote-staging-release.sh} \texttt{vX.Y.Z}.
    \item Commit staging overlay vào CD repo.
    \item Operator review diff và sync \texttt{yas-staging} thủ công.
\end{enumerate}

\lstinputlisting[
  caption={Operator sync staging sau khi release GitOps diff được tạo},
  basicstyle=\ttfamily\scriptsize
]{../evidence/19_staging_manual_sync_applied.txt}

\begin{figure}[H]
\centering
\reportimg{img/jenkins/release_v013_success.png}
\caption{Jenkins release tag \texttt{v0.1.3} promote thành công}
\end{figure}

\begin{figure}[H]
\centering
\reportimg{img/argocd/24_argocd_outofsync_diff.png}
\caption{ArgoCD staging có diff trước khi operator sync thủ công}
\end{figure}

\begin{table}[H]
\centering
\caption{Quá trình hardening staging release}
\begin{tabular}{|p{2.4cm}|p{4.5cm}|p{6.4cm}|}
\hline
\textbf{Lần chạy} & \textbf{Kết quả} & \textbf{Bài học kỹ thuật} \\
\hline
\texttt{v0.1.0} & GitOps gate chặn. & Tag đóng vào commit cũ làm lệch service scope giữa app repo và CD repo. \\
\hline
\texttt{v0.1.1} lần 1 & Jenkins fail do thiếu buildx. & Release promote cần agent có Docker buildx/imagetools. \\
\hline
\texttt{v0.1.1} lần 2 & Fail do mixed dev tags. & Không thể giả định mọi service trong dev cùng một SHA. \\
\hline
\texttt{v0.1.3} & Thành công. & Promote theo tag nguồn từng service đọc từ dev overlay; release tag là promote-only. \\
\hline
\end{tabular}
\end{table}

\subsection{Developer Preview}

Yêu cầu đề bài có job \texttt{developer\_build} để developer nhập branch cho service
muốn test. Nhóm tách job này sang CD repo để giữ main app pipeline đơn giản:

\begin{itemize}
    \item App repo feature branch chỉ build/push image commit SHA.
    \item CD repo \path{Jenkinsfile.developer-build} resolve branch sang SHA và kiểm tra image đã tồn tại trên Docker Hub.
    \item Job gọi \path{scripts/prepare-developer-preview.sh} \texttt{tax=<sha>} hoặc nhiều assignment khác.
    \item \texttt{developer} active, \texttt{dev} vẫn active, \texttt{staging} dormant để tránh quá tải.
    \item ArgoCD tự sync \texttt{yas-developer}; developer truy cập \url{http://storefront.developer.yas.local.com:30846/}.
\end{itemize}

\begin{figure}[H]
\centering
\reportimg{img/developer/32_developer_reachable.png}
\caption{Developer preview truy cập được qua NodePort 30846}
\end{figure}

\begin{figure}[H]
\centering
\reportimg{img/jenkins/developer_build_parameters.png}
\caption{Giao diện chọn tham số nhánh (Parameters) cho các service của Job developer\_build trên Jenkins}
\end{figure}

\begin{figure}[H]
\centering
\reportimg{img/argocd/developer_active_staging_dormant.png}
\caption{Giao diện ArgoCD hiển thị môi trường developer được kích hoạt (Active) và staging tạm tắt (Dormant)}
\end{figure}

\subsection{Teardown Developer}

Job \texttt{teardown\_developer} chạy từ CD repo, yêu cầu confirm
\texttt{developer-teardown}. Job reset overlay developer về tag \texttt{main},
đưa \texttt{developer} về dormant và khôi phục baseline \texttt{dev + staging active}.
Đây là cách đáp ứng yêu cầu xóa phần triển khai developer mà vẫn không dùng
\texttt{kubectl delete} trực tiếp.

\begin{figure}[H]
\centering
\reportimg{img/jenkins/teardown_developer_success.png}
\caption{Console log chạy thành công Job teardown\_developer khôi phục baseline môi trường}
\end{figure}

\subsection{Bằng Chứng Demo 3 Luồng}

Ba luồng demo đều đã chạy thật và có evidence tương ứng. Bảng dưới đây gom lại
đường đi, cơ chế sync và bằng chứng chính để người đọc không phải suy luận từ
các hình rời rạc ở những phần trước.

\begin{table}[H]
\centering
\caption{Tổng hợp evidence cho ba luồng CD demo}
\begin{tabular}{|p{2.8cm}|p{5.2cm}|p{5.8cm}|}
\hline
\textbf{Luồng} & \textbf{Cơ chế chạy} & \textbf{Evidence trong report} \\
\hline
Dev auto deploy & Push/merge vào \texttt{main} app repo, Jenkins build image commit-SHA, commit overlay \texttt{dev}, ArgoCD \texttt{yas-dev} auto-sync. & Hình Jenkins main build, Docker Hub commit SHA tag, ArgoCD apps overview và Storefront dev reachable. \\
\hline
Staging release & Git tag \texttt{v0.1.3}, Jenkins promote image có sẵn sang release tag, commit overlay \texttt{staging}, operator manual-sync. & Hình Jenkins release tag \texttt{v0.1.3}, ArgoCD staging diff/manual sync và release hardening table. \\
\hline
Developer preview/teardown & Job \texttt{developer\_build} bật môi trường \texttt{developer} theo branch/service; job \texttt{teardown\_developer} đưa developer về dormant và khôi phục \texttt{dev + staging}. & Hình developer reachable, ArgoCD developer synced, console teardown success. \\
\hline
\end{tabular}
\end{table}

Tóm tắt evidence trong \path{evidence/35_cd_demo_summary.md}:

\begin{itemize}
    \item Dev auto deploy: merge vào \texttt{main} tạo image commit-SHA, Jenkins commit \texttt{cd(lab2): update dev image tags [skip ci]} vào CD repo và ArgoCD tự sync \texttt{yas-dev}.
    \item Staging release: tag \texttt{v0.1.3} chạy Jenkins promote-only thành công; staging dùng release tag immutable và cần operator sync thủ công.
    \item Developer preview: job \texttt{developer\_build} bật \texttt{yas-developer} theo branch/service; job \texttt{teardown\_developer} đưa developer về dormant và khôi phục baseline \texttt{dev + staging}.
    \item Các lỗi release cũ \texttt{v0.1.0}/\texttt{v0.1.1} được giữ làm evidence hardening: GitOps gate, thiếu buildx và lỗi mixed dev tags.
\end{itemize}

\begin{figure}[H]
\centering
\reportimg{img/argocd/dev_staging_developer_apps_detail.png}
\caption{ArgoCD hiển thị các application chính của dev, staging và developer sau demo}
\end{figure}

\begin{figure}[H]
\centering
\reportimg{img/jenkins/teardown_developer_success.png}
\caption{Job teardown\_developer chạy thành công và trả môi trường developer về dormant}
\end{figure}
